\section{Opgaver}
\subsection{Opgave 1}
Lav et dokument med overskriften "LaTeX aften" og tilhørende brødtekst. I brødteksten skal sætningen "Der er da ikke så svært?" gentages som henholdsvis almindelig, kursiv, fed og understreget tekst. Der skal være linjeskift mellem de forskellige slags tekst. Lav til sidst et sideskift til næste opgave.

Din løsning burde minde om følgende:
\begin{figure}[H]
    \includegraphics[scale=1.3]{Billeder/Opgave 1.png}
    \label{fig:opg1}
\end{figure}

\subsection{Opgave 2}
Indsæt Pythagoras' læresætning direkte i noget tekst. Bagefter skal du finde din yndlings ligning og indsætte dem som en decideret \textit{equation}, hvor den bliver centreret.

Din løsning burde minde om følgende:
\begin{figure}[H]
    \includegraphics[scale=0.7]{Billeder/Opgave 2.png}
    \label{fig:opg2}
\end{figure}

\subsection{Opgave 3}
Lav både en bulletpoint liste og en nummereret liste.

\subsection{Opgave 4}
Indsæt et billede i dit dokument. Prøv at skaler billedet så det ser bedst muligt ud. (Herfra er Google din bedste ven!)

\subsection{Opgave 5}
Lav en tabel. Prøv først at lave den manuelt og herefter med en tablegenerator.

\subsection{Opgave 6}
Indsæt et flertal af ligninger i samme \textit{equation} environment. Prøv derefter at bruge \textit{align} environment til at få ligmed-tegnet til at stå samme sted.
\begin{figure}[H]
    \centering
    \includegraphics[scale=0.6]{Billeder/Opgave 6.png}
    \label{fig:opg6}
\end{figure}

\subsection{Opgave 7}
Indsæt noget kode (eksempelvis Java, Python, Matlab etc.) fra enten din sidste grupperegning eller fra nettet.